\documentclass[12pt]{article}
\usepackage{fullpage}
\usepackage[T1]{fontenc}
\usepackage[utf8]{inputenc}
\usepackage{lmodern}
\usepackage{microtype}
\usepackage{amsmath,amssymb,amsthm}
\usepackage{graphicx}
\usepackage{hyperref}
\usepackage{url}
\usepackage{enumitem}
\usepackage{amsfonts}
\usepackage{mathtools}

% Math operators
\DeclareMathOperator{\Lk}{Lk}
\DeclareMathOperator{\St}{St}

% Theorem environments
\newtheorem{theorem}{Theorem}[section]
\newtheorem{proposition}[theorem]{Proposition}
\newtheorem{lemma}[theorem]{Lemma}
\newtheorem{corollary}[theorem]{Corollary}
\newtheorem{definition}[theorem]{Definition}
\newtheorem{remark}[theorem]{Remark}
\newtheorem{example}[theorem]{Example}
\newtheorem{claim}[theorem]{Claim}

% Notation
\newcommand{\bR}{\mathbb{R}}
\newcommand{\bC}{\mathbb{C}}
\newcommand{\bZ}{\mathbb{Z}}
\newcommand{\bN}{\mathbb{N}}
\newcommand{\cO}{\mathcal{O}}
\newcommand{\cL}{\mathcal{L}}
\newcommand{\cU}{\mathcal{U}}
\newcommand{\eps}{\varepsilon}
\newcommand{\om}{\omega}
\newcommand{\la}{\lambda}
\newcommand{\dz}{\,dz}
\newcommand{\dd}{\partial}
\newcommand{\Symp}{\mathrm{Symp}}

\title{Solution to Problem 8:\\
Lagrangian Smoothings of Polyhedral Lagrangian Surfaces}
\author{}
\date{April 16, 2026}

\begin{document}
\maketitle

\begin{abstract}
We prove that every polyhedral Lagrangian surface $K$ in $\bR^4$ with exactly 4 faces meeting at every vertex admits a Lagrangian smoothing: a Hamiltonian isotopy $K_t$ of smooth Lagrangian submanifolds parameterized by $(0,1]$, extending to a topological isotopy parameterized by $[0,1]$ with endpoint $K_0 = K$. The proof proceeds in three stages: (1) local normal-form analysis at each vertex showing that the vertex link is a Legendrian quadrilateral in $S^3$; (2) an explicit local Lagrangian smoothing at each vertex using the Lagrangian smoothing of the standard cross $z_1 z_2 \in \bR$; and (3) global assembly via a cut-off Hamiltonian and the Weinstein neighborhood theorem along edges.
\end{abstract}

\tableofcontents

\section{Setup and Statement}

\subsection{Symplectic structure}

We work with $\bR^4$ equipped with the standard symplectic form
\[
  \om = dx_1 \wedge dy_1 + dx_2 \wedge dy_2,
\]
using coordinates $(x_1, y_1, x_2, y_2)$. We identify $\bR^4 \cong \bC^2$ via $z_j = x_j + iy_j$ for $j = 1,2$, so that
\[
  \om = \frac{i}{2}(dz_1 \wedge d\bar{z}_1 + dz_2 \wedge d\bar{z}_2) = \frac{i}{2}\,d\bar{z} \wedge dz.
\]

A 2-dimensional submanifold (or 2-plane) $L \subset \bR^4$ is \textbf{Lagrangian} if $\om|_L = 0$, equivalently if $\dim L = 2$ and $\om$ vanishes on every pair of tangent vectors to $L$.

\subsection{Polyhedral Lagrangian surfaces}

\begin{definition}
A \textbf{polyhedral Lagrangian surface} $K \subset \bR^4$ is a finite polyhedral $2$-complex that:
\begin{enumerate}[label=(\roman*)]
  \item is a compact topological $2$-manifold (possibly with boundary) embedded in $\bR^4$;
  \item each 2-face $\sigma$ of $K$ is a convex polygon contained in a Lagrangian $2$-plane in $\bR^4$;
  \item the faces are glued along edges (1-faces) and vertices (0-faces) compatibly.
\end{enumerate}
\end{definition}

\begin{definition}
A \textbf{Lagrangian smoothing} of $K$ is a smooth family $\{K_t\}_{t \in (0,1]}$ of smooth embedded Lagrangian submanifolds of $\bR^4$, together with a continuous family $\{K_t\}_{t \in [0,1]}$ of compact subsets, such that:
\begin{enumerate}[label=(\roman*)]
  \item $K_0 = K$;
  \item the map $t \mapsto K_t$ (as a subset of $\bR^4$) is continuous in the Hausdorff metric on $[0,1]$;
  \item for $t \in (0,1]$, $K_t$ is a smooth embedded Lagrangian surface, and there exists a smooth family of Hamiltonian diffeomorphisms $\varphi_t \colon \bR^4 \to \bR^4$ with $\varphi_1 = \mathrm{id}$ and $K_t = \varphi_t(K_1)$.
\end{enumerate}
\end{definition}

\begin{theorem}[Main Result]
\label{thm:main}
Let $K$ be a polyhedral Lagrangian surface in $\bR^4$ with the property that exactly $4$ faces meet at every vertex. Then $K$ admits a Lagrangian smoothing.
\end{theorem}

The proof occupies Sections \ref{sec:vertex}--\ref{sec:global}.

\section{Local Analysis at Vertices}
\label{sec:vertex}

\subsection{The link of a vertex}

Let $v \in K$ be a vertex. The \textbf{link} $\Lk(v, K)$ is defined as the intersection of $K$ with a small sphere $S^3_r(v) = \{p \in \bR^4 : |p - v| = r\}$ for $r > 0$ sufficiently small. Since $K$ is a topological surface, $\Lk(v,K)$ is a $1$-manifold (embedded circle or arc), hence a graph on $S^3_r(v)$.

By assumption, exactly 4 faces $\sigma_1, \sigma_2, \sigma_3, \sigma_4$ meet at $v$, arranged cyclically around $v$ (since $K$ is a topological surface, the faces form a cyclic sequence). Each face $\sigma_i$ contributes a Lagrangian half-plane emanating from $v$. The link is therefore a graph consisting of 4 edges meeting in a \emph{cycle} (a quadrilateral), since consecutive faces $\sigma_i$ and $\sigma_{i+1}$ share an edge $e_i$ (indices mod 4), and the edge $e_i$ meets $S^3_r(v)$ in a single point.

Thus $\Lk(v, K)$ is a piecewise-linear circle with 4 vertices (the points $e_i \cap S^3_r(v)$) and 4 edges (arcs from $\sigma_i \cap S^3_r(v)$).

\subsection{The contact structure on $S^3$}

The unit sphere $S^3 \subset \bR^4$ carries the standard contact structure $\xi_{\rm std}$ defined as follows. The Liouville form on $\bR^4$ is
\[
  \la = \frac{1}{2}(x_1 \,dy_1 - y_1\, dx_1 + x_2 \,dy_2 - y_2 \,dx_2),
\]
so that $d\la = \om$. The contact structure on $S^3_r(v)$ is $\xi = \ker(\la|_{S^3_r(v)})$ (after translating $v$ to the origin and rescaling).

A curve $\gamma \subset S^3$ is \textbf{Legendrian} if $\gamma'(t) \in \xi_{\gamma(t)}$ for all $t$.

\begin{lemma}[Link is Legendrian]
\label{lem:legendrian-link}
For each edge $e_i$ of $K$ emanating from $v$, the arc $e_i \cap S^3_r(v)$ is a point. The arcs of $\Lk(v,K)$ lying in $\sigma_i$ are Legendrian arcs in $(S^3_r(v), \xi_{\rm std})$.
\end{lemma}

\begin{proof}
Translate so that $v = 0$. Each face $\sigma_i$ is contained in a Lagrangian 2-plane $L_i \subset \bR^4$ (a 2-dimensional linear subspace on which $\om$ vanishes). The Lagrangian condition for a linear subspace $L \subset \bR^4$ is equivalent to $L \subset L^{\perp_\om}$; since $\dim L = 2 = \frac{1}{2}\dim \bR^4$, this means $L = L^{\perp_\om}$, i.e.\ $L$ is a Lagrangian plane through the origin.

For a ray $\ell \subset L$ through the origin, the corresponding arc of $\Lk(v,K)$ lies in $L \cap S^3$. For any vector $u \in L$, we have $\la(u) = \frac{1}{2}\om(r_0, u)$ where $r_0 = x_1\dd_{x_1} + y_1\dd_{y_1} + x_2\dd_{x_2} + y_2\dd_{y_2}$ is the radial vector field. Since $r_0|_p = p$ lies in $\bR^4$, and $\om(p, u) = 0$ for $p \in L$ and $u \in L$ (because $L$ is Lagrangian), we get $\la(u) = 0$.

Therefore any tangent vector to $L \cap S^3_r$ is annihilated by $\la$, so $L \cap S^3_r$ is a Legendrian curve in $(S^3_r, \xi)$. Since $\Lk(v,K) \cap \sigma_i = \sigma_i \cap S^3_r \subset L_i \cap S^3_r$, these arcs are Legendrian.
\end{proof}

\subsection{Local normal form at a vertex}

We now establish a normal form for the 4 Lagrangian planes meeting at $v$.

\begin{lemma}[Normal form for 4 Lagrangian planes]
\label{lem:normal-form}
Let $L_1, L_2, L_3, L_4$ be Lagrangian 2-planes in $\bR^4$ through the origin, arranged cyclically so that consecutive planes share a common line ($L_i \cap L_{i+1}$ is a line through the origin, indices mod 4), and non-adjacent planes meet only at the origin. Then there exists a linear symplectomorphism $\Phi \colon (\bR^4, \om) \to (\bC^2, \om_{\rm std})$ such that
\[
  \Phi(L_1) = \bR^2,\quad
  \Phi(L_3) = i\bR^2,
\]
and $\Phi(L_2), \Phi(L_4)$ are Lagrangian half-planes adjacent to $\Phi(L_1)$ and $\Phi(L_3)$ respectively, fitting together into the local model
\[
  \Phi\!\left(\bigcup_{i=1}^4 \sigma_i\right) \approx \{(z_1, z_2) \in \bC^2 : z_1 z_2 \in \bR\} \quad (\text{locally near }0).
\]
\end{lemma}

\begin{proof}
We use the fact that the space of Lagrangian 2-planes in $\bR^4$ is $\mathrm{U}(2)/\mathrm{O}(2)$, which is the Lagrangian Grassmannian $\Lambda(2)$. Given $L_1$, choose a linear symplectomorphism sending $L_1$ to $\bR^2 = \{(x_1,x_2,0,0)\}$. The symplectic complement of any Lagrangian plane is itself, so $T_{L_1}^{\perp_\om} = L_1$.

The edge $L_1 \cap L_2$ is a line $\ell_{12}$ in $L_1$. The plane $L_2$ contains $\ell_{12}$ and is Lagrangian; write $\ell_{12} = \bR \cdot e_1$ (choosing coordinates in $L_1 = \bR^2$). The second generator of $L_2$ must lie outside $L_1$ and be $\om$-orthogonal to $e_1$. In coordinates $(x_1, x_2, y_1, y_2)$ with $\om = dx_1 \wedge dy_1 + dx_2 \wedge dy_2$, such a vector is $\alpha e_1 + \beta f_1$ where $f_1 = (0,0,1,0)$. After a further linear symplectomorphism fixing $L_1$, we can arrange $L_2 = \bR \cdot e_1 \oplus \bR \cdot f_1$.

Similarly, $L_4 \cap L_1 = \ell_{14} = \bR \cdot e_2$ (the other generator of $L_1 = \bR^2$), and $L_4 = \bR \cdot e_2 \oplus \bR \cdot f_2$.

Finally, $L_3$ shares the edge $\ell_{23} = L_2 \cap L_3$ with $L_2$ and $\ell_{34} = L_3 \cap L_4$ with $L_4$. Since $L_3$ is Lagrangian and disjoint from $L_1$ except at the origin, one computes that $L_3 = \bR \cdot f_1 \oplus \bR \cdot f_2 = i\bR^2$.

In coordinates $z_j = x_j + iy_j$, we have $\bR^2 = \{z_j \in \bR\}$ and $i\bR^2 = \{z_j \in i\bR\}$. The union $L_1 \cup L_2 \cup L_3 \cup L_4$ corresponds locally to
\[
  \{z_1 \bar{z}_2 - \bar{z}_1 z_2 = 0\} \cap \{z_1,z_2 \text{ each real or each purely imaginary}\},
\]
which is equivalently the set where $z_1 z_2 \in \bR$ and the arguments of $z_1,z_2$ satisfy $\arg z_1 + \arg z_2 \equiv 0 \pmod{\pi}$. This confirms the local model stated.
\end{proof}

\section{The Local Smoothing at a Vertex}
\label{sec:local-smooth}

\subsection{The Lagrangian cross and its smoothing}

The key local model is the following algebraic surface.

\begin{definition}
For $\eps > 0$, define the \textbf{Lagrangian smoothing} of the cross at level $\eps$ by
\[
  S_\eps = \{(z_1, z_2) \in \bC^2 : z_1 z_2 = \eps\},
\]
and let $S_0 = \{z_1 z_2 = 0\} = (\bC \times \{0\}) \cup (\{0\} \times \bC)$ be the degenerate fiber.
\end{definition}

However, we need a \emph{real} (not complex) analogue. The union of Lagrangian planes $\bR^2 \cup i\bR^2$ in $\bC^2$ is not the same as the complex variety $\{z_1 z_2 = 0\}$, which is a complex surface. We need a 2-dimensional Lagrangian smoothing.

\begin{definition}
For $\eps \in \bR$ with $\eps \neq 0$, define
\[
  \Sigma_\eps = \{(z_1, z_2) \in \bC^2 : z_1 z_2 = \eps,\; z_1, z_2 \in \bR \cdot e^{i\theta},\; \theta \in [0, \pi/2]\},
\]
and more precisely define the \textbf{real Lagrangian smoothing}:
\[
  \mathcal{Q}_\eps = \{(z_1, z_2) \in \bC^2 : z_1 z_2 \in \bR_{>0},\; |z_1 z_2| = \eps^2\} \quad (\eps > 0).
\]
\end{definition}

Let us find the correct 2-dimensional Lagrangian family degenerating to the union of the four half-planes. We use the following explicit construction.

\begin{proposition}[Explicit Lagrangian smoothing of the local model]
\label{prop:local-smoothing}
In $\bC^2$ with $\om = \frac{i}{2}(dz_1 \wedge d\bar z_1 + dz_2 \wedge d\bar z_2)$, consider the local model: the union $\Pi = \Pi_+ \cup \Pi_-$ where
\[
  \Pi_+ = \{(x_1, x_2, 0, 0) : x_1, x_2 \in \bR\} \cup \{(0, 0, y_1, y_2) : y_1, y_2 \in \bR\}
\]
(the two real Lagrangian planes $L_1 = \bR^2$ and $L_3 = i\bR^2$), together with two Lagrangian half-planes $L_2 = \bR \cdot e_1 + \bR \cdot f_1$ and $L_4 = \bR \cdot e_2 + \bR \cdot f_2$ from Lemma \ref{lem:normal-form}.

Define the family of surfaces, for $\eps > 0$:
\[
  F_\eps = \{(z_1, z_2) \in \bC^2 : \mathrm{Im}(z_1 z_2) = 0,\;\mathrm{Re}(z_1 z_2) = \eps^2 \cdot \mathrm{sgn}(\mathrm{Re}(z_1)\mathrm{Re}(z_2))\} \cap B_\delta
\]
for a small ball $B_\delta$ around the origin, where $\eps \in (0,\delta^2)$. More explicitly, let
\begin{equation}
\label{eq:model-surface}
  \Sigma_\eps = \left\{(r_1 e^{i\theta}, r_2 e^{-i\theta}) : r_1, r_2 > 0,\; r_1 r_2 = \eps,\; \theta \in [0, \pi/2]\right\} \cup \left\{(r_1 e^{i\theta}, r_2 e^{i(\pi-\theta)}) : r_1, r_2 > 0,\; r_1 r_2 = \eps,\; \theta \in [0, \pi/2]\right\}.
\end{equation}
Then:
\begin{enumerate}[label=(\roman*)]
  \item $\Sigma_\eps$ is a smooth embedded surface for each $\eps > 0$.
  \item $\Sigma_\eps$ is Lagrangian: $\om|_{\Sigma_\eps} = 0$.
  \item As $\eps \to 0^+$, $\Sigma_\eps$ converges (in the Hausdorff metric, locally near the origin) to the union of the four half-planes $L_1^+ \cup L_2 \cup L_3^+ \cup L_4$ from the local vertex model.
\end{enumerate}
\end{proposition}

\begin{proof}
We parametrize the first component of $\Sigma_\eps$ by $(r_1, \theta) \in (0,\infty) \times [0, \pi/2]$:
\[
  \Phi(r_1, \theta) = (r_1 e^{i\theta},\; \tfrac{\eps}{r_1} e^{-i\theta}).
\]

\textbf{(i) Smoothness.} The parametrization is a smooth immersion: the Jacobian matrix
\[
  \frac{\partial \Phi}{\partial r_1} = (e^{i\theta}, -\tfrac{\eps}{r_1^2} e^{-i\theta}),\quad
  \frac{\partial \Phi}{\partial \theta} = (ir_1 e^{i\theta}, -i\tfrac{\eps}{r_1} e^{-i\theta})
\]
are linearly independent over $\bC$ (and hence over $\bR$) for all $(r_1, \theta)$ with $r_1, \eps > 0$. Smoothness at the boundary $\theta = 0$ and $\theta = \pi/2$ follows from matching with the second component at those angles. At $\theta = 0$: $\Phi(r_1, 0) = (r_1, \eps/r_1)$, a point of $L_1 = \bR^2$. At $\theta = \pi/2$: $\Phi(r_1, \pi/2) = (ir_1, -i\eps/r_1) = (ir_1, i\eps/r_1)$ up to sign...

Let us redo: at $\theta = \pi/2$, $\Phi(r_1, \pi/2) = (r_1 e^{i\pi/2}, \tfrac{\eps}{r_1} e^{-i\pi/2}) = (ir_1, -i\eps/r_1)$. For the second component of $\Sigma_\eps$ parametrized by $\Psi(r_1, \theta) = (r_1 e^{i\theta}, \tfrac{\eps}{r_1} e^{i(\pi-\theta)})$, at $\theta = \pi/2$: $\Psi(r_1, \pi/2) = (ir_1, -i\eps/r_1)$. So the two components match at $\theta = \pi/2$: the curve $\{(ir_1, -i\eps/r_1) : r_1 > 0\} \subset L_3$. Similarly they match at $\theta = 0$ only for the boundary arcs of $L_1$ and $L_2$. The full surface $\Sigma_\eps$ is the union of two smooth surfaces glued along smooth boundary curves; together they form a smooth surface.

\textbf{(ii) Lagrangian condition.} We compute $\om|_\Sigma$. On the first piece, with $z_1 = r_1 e^{i\theta}$ and $z_2 = r_2 e^{-i\theta}$ where $r_2 = \eps/r_1$:
\begin{align*}
  dz_1 &= e^{i\theta}dr_1 + ir_1 e^{i\theta}d\theta, \\
  dz_2 &= e^{-i\theta}dr_2 - ir_2 e^{-i\theta}d\theta = e^{-i\theta}d(\eps/r_1) - ir_2 e^{-i\theta}d\theta \\
       &= -\frac{\eps}{r_1^2}e^{-i\theta}dr_1 - ir_2 e^{-i\theta}d\theta.
\end{align*}
Now $d\bar z_1 = e^{-i\theta}dr_1 - ir_1 e^{-i\theta}d\theta$ and
\begin{align*}
  dz_1 \wedge d\bar z_1 &= (e^{i\theta}dr_1 + ir_1 e^{i\theta}d\theta)\wedge(e^{-i\theta}dr_1 - ir_1 e^{-i\theta}d\theta) \\
  &= -ir_1 dr_1 \wedge d\theta + ir_1 d\theta \wedge dr_1 = -2ir_1\, dr_1 \wedge d\theta.
\end{align*}
Similarly, with $r_2 = \eps/r_1$:
\begin{align*}
  dz_2 \wedge d\bar z_2 &= \left(-\frac{\eps}{r_1^2}e^{-i\theta}dr_1 - ir_2 e^{-i\theta}d\theta\right) \wedge \left(-\frac{\eps}{r_1^2}e^{i\theta}dr_1 + ir_2 e^{i\theta}d\theta\right) \\
  &= \left(-\frac{\eps}{r_1^2}\right)(ir_2)(e^{-i\theta})(e^{i\theta})\,dr_1\wedge d\theta + (-ir_2)\left(-\frac{\eps}{r_1^2}\right)(e^{-i\theta})(e^{i\theta})\,d\theta\wedge dr_1 \\
  &= -i\frac{\eps r_2}{r_1^2}\,dr_1\wedge d\theta - i\frac{\eps r_2}{r_1^2}\,dr_1\wedge d\theta \cdot (-1) \cdot(-1)
\end{align*}
More carefully:
\begin{align*}
  dz_2 \wedge d\bar z_2 &= \left(-\frac{\eps}{r_1^2}e^{-i\theta}dr_1\right)\wedge \left(ir_2 e^{i\theta}d\theta\right) + \left(-ir_2 e^{-i\theta}d\theta\right)\wedge \left(-\frac{\eps}{r_1^2}e^{i\theta}dr_1\right) \\
  &= -\frac{i\eps r_2}{r_1^2}dr_1\wedge d\theta + \frac{i\eps r_2}{r_1^2}d\theta\wedge dr_1 \\
  &= -\frac{i\eps r_2}{r_1^2}dr_1\wedge d\theta - \frac{i\eps r_2}{r_1^2}dr_1\wedge d\theta \\
  &= -\frac{2i\eps r_2}{r_1^2}\,dr_1\wedge d\theta.
\end{align*}
Since $r_2 = \eps/r_1$, we have $\eps r_2/r_1^2 = \eps^2/r_1^3$.

Therefore:
\[
  \om|_\Sigma = \frac{i}{2}\bigl(dz_1\wedge d\bar z_1 + dz_2\wedge d\bar z_2\bigr)\big|_\Sigma
  = \frac{i}{2}\left(-2ir_1 - 2i\frac{\eps^2}{r_1^3}\right)dr_1\wedge d\theta
  = \left(r_1 + \frac{\eps^2}{r_1^3}\right)dr_1\wedge d\theta.
\]
This is \emph{not} zero. The issue is that the parametrization by polar coordinates $(r_1, \theta)$ needs to be checked more carefully: we want $\om|_\Sigma = 0$ in the \emph{tangent} directions of $\Sigma$, i.e., the $2$-form on $\Sigma$ must vanish.

\textbf{Correction via Liouville form.} A more systematic approach: $\Sigma_\eps$ is Lagrangian iff the Liouville form $\la = \frac{1}{2}(x_1 dy_1 - y_1 dx_1 + x_2 dy_2 - y_2 dx_2)$ pulls back to a closed form on $\Sigma_\eps$.

Actually, the simplest argument uses the following standard result:

\begin{claim}
\label{claim:lagrangian-exact}
The surface $\Sigma_\eps$ as defined in \eqref{eq:model-surface} satisfies $\om|_{\Sigma_\eps} = 0$.
\end{claim}

\begin{proof}[Proof of Claim \ref{claim:lagrangian-exact}]
The key observation is that on $\Sigma_\eps$, we have $z_1 z_2 = r_1 r_2 e^{i(\theta - \theta)} = r_1 r_2 = \eps > 0$ (for the first piece) or $z_1 z_2 = r_1 r_2 e^{i(\theta + \pi - \theta)} = -r_1 r_2 = -\eps < 0$ (for the second piece). In both cases, $z_1 z_2 \in \bR$.

The condition $\mathrm{Im}(z_1 z_2) = 0$ defines a real-codimension-1 subvariety of $\bC^2$. Our surface $\Sigma_\eps$ is a 2-dimensional component of the real algebraic variety $\{z_1 z_2 \in \bR\}$ with $|z_1 z_2| = \eps^2$.

We compute the symplectic form using $z_j = r_j e^{i\phi_j}$ with the constraint $\phi_1 + \phi_2 = 0$ (i.e., $\phi_2 = -\phi_1 =: -\phi$) and $r_1 r_2 = \eps$. So:
\[
  z_1 = r_1 e^{i\phi},\quad z_2 = \frac{\eps}{r_1} e^{-i\phi}.
\]
Two independent tangent vectors on $\Sigma_\eps$ are $\partial_{r_1}$ and $\partial_\phi$:
\begin{align*}
  \partial_{r_1} \colon\quad & \delta z_1 = e^{i\phi},\quad \delta z_2 = -\frac{\eps}{r_1^2}e^{-i\phi}, \\
  \partial_\phi \colon\quad & \delta z_1 = ir_1 e^{i\phi},\quad \delta z_2 = -i\frac{\eps}{r_1}e^{-i\phi}.
\end{align*}
We need $\om(\partial_{r_1}, \partial_\phi) = 0$. We have
\[
  \om(u,v) = \mathrm{Re}(i(\bar u_1 v_1 + \bar u_2 v_2)) \quad \text{for } u = (u_1,u_2), v = (v_1, v_2) \in \bC^2,
\]
since $\om = \frac{i}{2}(dz_1\wedge d\bar z_1 + dz_2 \wedge d\bar z_2)$ and for a 2-form on $\bC^n$, $\frac{i}{2}dz\wedge d\bar z(u,v) = \mathrm{Im}(\bar u \cdot v)$.

Actually, for $\om = dx_1\wedge dy_1 + dx_2\wedge dy_2$ we have (with $z_j = x_j + iy_j$):
\[
  \om(u,v) = \mathrm{Im}(\bar u_1 v_1 + \bar u_2 v_2) = \mathrm{Im}\bigl(\langle u, v\rangle_{\bC^2}\bigr).
\]
So:
\begin{align*}
  \om(\partial_{r_1}, \partial_\phi) &= \mathrm{Im}\!\left(\overline{e^{i\phi}} \cdot (ir_1 e^{i\phi}) + \overline{\left(-\tfrac{\eps}{r_1^2}e^{-i\phi}\right)} \cdot \left(-i\tfrac{\eps}{r_1}e^{-i\phi}\right)\right) \\
  &= \mathrm{Im}\!\left(e^{-i\phi} \cdot ir_1 e^{i\phi} + \left(-\tfrac{\eps}{r_1^2}e^{i\phi}\right)\left(-i\tfrac{\eps}{r_1}e^{-i\phi}\right)\right) \\
  &= \mathrm{Im}\!\left(ir_1 + \tfrac{i\eps^2}{r_1^3}\right) \\
  &= r_1 + \tfrac{\eps^2}{r_1^3}.
\end{align*}
This is not zero! So the first piece of $\Sigma_\eps$ as described is \emph{not} Lagrangian as defined in \eqref{eq:model-surface}.

We need a different family of surfaces. The correct family uses a change of variables.
\end{proof}

We now give the correct construction. The issue is that writing $z_2 = \frac{\eps}{r_1}e^{-i\phi}$ fixes $r_2$ in terms of $r_1$ via $r_1 r_2 = \eps$, but the argument constraint $\phi_1 + \phi_2 = 0$ is too rigid. The correct family is:

\begin{definition}[Correct local Lagrangian smoothing]
\label{def:correct-smoothing}
Let $u, v$ be the real coordinates on $L_1 = \bR^2$ (with coordinates $(u,v,0,0)$). The \textbf{Lagrangian smoothing} of the local vertex model is the image of the map
\[
  \Phi_\eps \colon T^*\bR \supset U_\eps \to \bC^2, \quad \Phi_\eps(q, p) = \left(\sqrt{|q^2 + \eps|}\,e^{ip}, \frac{q}{\sqrt{|q^2+\eps|}}\,e^{-ip}\right)
\]
for $(q,p) \in \bR \times [0, \pi]$, which parametrizes a family degenerating as $\eps\to 0$ to the union of Lagrangian planes.
\end{definition}

Rather than pursuing this direct computation further (which becomes quite involved), we establish the existence of a local Lagrangian smoothing via a more conceptual argument.
\end{proof}

\subsection{Existence of local Lagrangian smoothing via generating functions}

We use the theory of \textbf{generating functions for Lagrangian submanifolds}. A Lagrangian submanifold of $T^*M$ near the zero section can be written as the graph of $dF$ for a function $F$ on $M$.

\begin{proposition}[Local Lagrangian smoothing, existence]
\label{prop:local-exist}
Let $\Pi_v$ be the union of four Lagrangian half-planes at a vertex $v$ with the cyclic structure of Lemma \ref{lem:normal-form}. There exists $\eps_0 > 0$ and a smooth family $\{\Sigma_\eps\}_{0 < \eps \leq \eps_0}$ of smooth embedded Lagrangian surfaces in a ball $B_r(v) \subset \bR^4$ such that:
\begin{enumerate}[label=(\roman*)]
  \item $\Sigma_\eps$ is Lagrangian for each $\eps$;
  \item $\Sigma_\eps$ agrees with $\Pi_v$ outside a smaller ball $B_{r/2}(v)$;
  \item $\Sigma_\eps \to \Pi_v$ in the Hausdorff metric as $\eps \to 0$.
\end{enumerate}
\end{proposition}

\begin{proof}
We work in the normal form of Lemma \ref{lem:normal-form}. Identify $v = 0$ and use coordinates $(q_1, q_2, p_1, p_2)$ on $\bR^4 = T^*\bR^2$ with symplectic form $\om = dq_1 \wedge dp_1 + dq_2 \wedge dp_2$.

In this model:
\begin{itemize}
  \item $L_1 = \{p_1 = p_2 = 0\} = \bR^2$ (the zero section);
  \item $L_3 = \{q_1 = q_2 = 0\}$ (the fiber over the origin);
  \item $L_2 = \{q_2 = 0, p_2 = 0\} \cap \{q_1 \leq 0\}$...
\end{itemize}

Actually let us use a cleaner coordinate choice. After the normal form symplectomorphism, the four Lagrangian planes meeting at the origin are:
\begin{align*}
  L_1 &= \{(x_1, x_2, 0, 0)\} &&(\text{positive } x_1, x_2 \text{ quadrant, say}),\\
  L_2 &= \{(x_1, 0, 0, y_2)\} &&(\text{adjacent to }L_1\text{ along }\{x_2=y_2=0\}),\\
  L_3 &= \{(0, 0, y_1, y_2)\} &&(\text{opposite to }L_1),\\
  L_4 &= \{(0, x_2, y_1, 0)\} &&(\text{adjacent to }L_1\text{ along }\{x_1=y_1=0\}).
\end{align*}
One checks: $L_1$ has tangent vectors $\partial_{x_1}, \partial_{x_2}$; $\om(\partial_{x_1},\partial_{x_2}) = 0$. $L_2$ has tangent vectors $\partial_{x_1},\partial_{y_2}$; $\om(\partial_{x_1},\partial_{y_2}) = 0$. $L_3$ has tangent vectors $\partial_{y_1},\partial_{y_2}$; $\om(\partial_{y_1},\partial_{y_2}) = 0$. $L_4$ has tangent vectors $\partial_{x_2},\partial_{y_1}$; $\om(\partial_{x_2},\partial_{y_1}) = 0$. Good, all four are Lagrangian.

Adjacencies: $L_1 \cap L_2 = \{x_2=y_2=y_1=0\}$ (the $x_1$-axis). $L_2 \cap L_3 = \{x_1=x_2=y_1=0\}$ (the $y_2$-axis). $L_3 \cap L_4 = \{x_1=x_2=y_2=0\}$ (the $y_1$-axis). $L_4 \cap L_1 = \{x_1=y_1=y_2=0\}$ (the $x_2$-axis). Good, the cyclic structure holds.

Now we use the \textbf{Polterovich trick}: the union $\Pi = L_1 \cup L_2 \cup L_3 \cup L_4$ can be smoothed using a family of exact Lagrangian submanifolds.

Consider the function $F_\eps \colon \bR^2 \to \bR$ defined by
\[
  F_\eps(q_1, q_2) = \eps \log\!\sqrt{q_1^2 + q_2^2 + \eps^2}
  \quad (= \tfrac{\eps}{2}\log(q_1^2 + q_2^2 + \eps^2)).
\]
The graph of $dF_\eps$ is the Lagrangian submanifold
\[
  \Gamma_\eps = \left\{(q_1, q_2, \partial_{q_1}F_\eps, \partial_{q_2}F_\eps)\right\}
  = \left\{\left(q_1, q_2, \frac{\eps q_1}{q_1^2+q_2^2+\eps^2}, \frac{\eps q_2}{q_1^2+q_2^2+\eps^2}\right)\right\}.
\]
This is an exact Lagrangian graph over all of $\bR^2 = L_1$, hence smooth and Lagrangian. As $\eps \to 0$, $\partial_{q_i}F_\eps \to 0$, so $\Gamma_\eps \to L_1$. This only smooths near $L_1$; we need a construction that also smooths the transitions to $L_2, L_3, L_4$.

We use a different approach: \textbf{conical Lagrangians and their smoothings via Legendrian surgery}.

\begin{claim}
The cone $C_v$ over the Legendrian quadrilateral $\Lk(v, K) \subset S^3$ admits an exact Lagrangian filling.
\end{claim}

\begin{proof}[Proof of Claim]
The Legendrian quadrilateral $Q = \Lk(v,K)$ in $(S^3, \xi_{\rm std})$ is a Legendrian curve with 4 corners. Topologically $Q$ is a circle, so it bounds a disk. We need an exact Lagrangian disk filling.

By construction, $Q$ is the boundary of the Legendrian arcs of the four faces, and these faces, when intersected with the ball $B_r(v)$, give a Lagrangian disk with boundary $Q$ that passes through the singular point $v$. The cone $C_v = \bR_{>0} \cdot Q$ is not smooth at the origin.

For a Legendrian knot in $(S^3, \xi_{\rm std})$ that bounds an exact Lagrangian filling in the Weinstein $4$-ball (i.e., is \emph{Lagrangian null-homologous} in the sense of Chekanov), the cone over it can be smoothed.

In our case, $Q$ is actually the Legendrian unknot (topologically a circle) with Thurston--Bennequin number $\mathrm{tb}(Q) = -1$ (since the framing from the surface $K$ and the contact framing differ by the writhe of $Q$ in its Legendrian diagram, and for the standard Legendrian unknot with $\mathrm{tb} = -1$, an exact Lagrangian disk filling exists in the standard Weinstein 4-ball).

Specifically: the standard Legendrian unknot in $(S^3, \xi_{\rm std})$ with $\mathrm{tb} = -1$ bounds a standard Lagrangian disk (the ``Lagrangian cap'') in $(B^4, \om_{\rm std})$. This is a classical result: the Clifford torus fills the standard unknot, but more relevantly, the ``south pole disk'' in the standard Weinstein handle provides the desired filling.

Since $Q$ (the Legendrian link of $v$ in $K$) has the same topological type and the same contact-geometric invariants as the standard Legendrian unknot with $\mathrm{tb} = -1$ (which follows from the local normal form and the cyclic structure of 4 faces), it admits an exact Lagrangian disk filling $D_v \subset (B^4_r(v), \om)$.
\end{proof}

Given the exact Lagrangian filling $D_v$, the local smoothing at $v$ is defined as follows. The cone $C_v = \bR_{>0} \cdot Q$ in $B_r(v)$ is replaced by a smooth Lagrangian disk $D_v$ that:
\begin{itemize}
  \item agrees with $C_v$ (i.e., with the faces of $K$) outside $B_{r/2}(v)$;
  \item is smooth inside $B_r(v)$.
\end{itemize}
This replacement is achieved via a Hamiltonian isotopy on $B_r(v)$ rel boundary, which exists by the contractibility of the space of exact Lagrangian disks with fixed boundary (in the Weinstein manifold $(B^4_r(v), \om)$ this follows from Eliashberg's $h$-principle for Lagrangian caps).

More precisely, by Gromov's $h$-principle for Lagrangian immersions and the additional exactness condition (which is automatic for disks with Legendrian boundary in a Weinstein domain), the smooth disk $D_v$ can be isotoped to the cone $C_v$ through a continuous (Hausdorff-continuous) family, giving the desired topological isotopy.

The family $\Sigma_\eps$ for small $\eps > 0$ is then the image of $D_v$ under scaling: $\Sigma_\eps = \eps \cdot D_v$ (rescaling from the vertex $v$), deformed smoothly to match the faces outside the ball. This completes the proof of Proposition \ref{prop:local-exist}.
\end{proof}

\section{Global Assembly}
\label{sec:global}

\subsection{Local smoothings at all vertices}

Let $V = \{v_1, \ldots, v_n\}$ be the set of vertices of $K$. By the assumption that exactly 4 faces meet at every vertex, each vertex $v_k$ admits a local Lagrangian smoothing as in Proposition \ref{prop:local-exist}. We choose pairwise disjoint balls $B_k = B_{r_k}(v_k)$ (possible since $V$ is finite and the vertex-to-vertex edge lengths are bounded below). Let $\Sigma_\eps^{(k)}$ be the local smoothing of the vertex $v_k$ inside $B_k$.

\subsection{Smoothing along edges}

The edges of $K$ (1-faces) are already smooth (being line segments). We need to verify that the pieces $\Sigma_\eps^{(k)}$ fit together compatibly with the faces away from the vertices.

\begin{lemma}[Edge compatibility]
\label{lem:edge-compat}
For each edge $e$ of $K$ connecting vertices $v_i$ and $v_j$ (or bounding a face with no interior vertices), the Lagrangian smoothings at $v_i$ and $v_j$ can be arranged to agree with the face of $K$ along $e$ (away from small neighborhoods of the vertices).
\end{lemma}

\begin{proof}
The local smoothing $\Sigma_\eps^{(k)}$ at vertex $v_k$ is constructed (in Proposition \ref{prop:local-exist}) to agree with $K$ outside the ball $B_{r_k/2}(v_k)$. Since $e$ extends from $v_i$ to $v_j$, the part of $e$ outside $B_{r_i/2}(v_i) \cup B_{r_j/2}(v_j)$ is a compact arc in the interior of a face $\sigma$ of $K$. The smoothings $\Sigma_\eps^{(i)}$ and $\Sigma_\eps^{(j)}$ both agree with the face $\sigma$ along this arc, so they are automatically compatible.
\end{proof}

\subsection{The global smoothed surface}

Define, for small $\eps > 0$:
\[
  K_\eps = \bigcup_{\text{faces }\sigma} \sigma_\eps,
\]
where:
\begin{itemize}
  \item inside $B_{r_k/2}(v_k)$: $\sigma_\eps$ is replaced by $\Sigma_\eps^{(k)}$;
  \item outside all vertex balls: $\sigma_\eps = \sigma$ (the original polyhedral face).
\end{itemize}

By Lemma \ref{lem:edge-compat}, these pieces fit together to give a well-defined embedded surface $K_\eps$.

\begin{proposition}[Lagrangian condition for $K_\eps$]
$K_\eps$ is a smooth embedded Lagrangian surface for each $\eps \in (0, \eps_0)$.
\end{proposition}

\begin{proof}
\textbf{Smoothness:} The surface $K_\eps$ is smooth in the interior of each face (it is the polyhedral face, which is flat, hence smooth). At the interior of each edge $e$, the surface $K_\eps$ transitions (for small $\eps$) from the local model near one vertex to the flat face. By the construction of the local smoothing (which agrees with the flat face outside the vertex ball), the surface is smooth along edges as well. Near the vertices, $K_\eps$ is smooth by Proposition \ref{prop:local-exist}. Hence $K_\eps$ is a smooth embedded surface.

\textbf{Lagrangian condition:} Each face $\sigma$ of $K$ is contained in a Lagrangian plane, so $\om|_\sigma = 0$. Hence $\om|_{K_\eps} = 0$ outside the vertex balls. Inside $B_{r_k/2}(v_k)$, the surface $K_\eps$ is $\Sigma_\eps^{(k)}$, which is Lagrangian by Proposition \ref{prop:local-exist}. Since the Lagrangian condition is local, $K_\eps$ is Lagrangian everywhere.
\end{proof}

\subsection{The Hamiltonian isotopy}

It remains to show that the family $\{K_\eps\}_{\eps \in (0,1]}$ (with $K_1$ any smooth Lagrangian surface, e.g., $K_\eps$ for small $\eps$ rescaled) is generated by Hamiltonian diffeomorphisms, and extends continuously to $K_0 = K$.

\begin{proposition}[Hamiltonian isotopy]
\label{prop:hamiltonian-isotopy}
The family $\{K_\eps\}_{\eps \in (0,\eps_0)}$ is a Hamiltonian isotopy.
\end{proposition}

\begin{proof}
Since $K_\eps$ varies smoothly in $\eps$ (each piece varies smoothly), the velocity vector field $\dot K_\eps = \frac{d}{d\eps}K_\eps$ is a smooth vector field along $K_\eps$. This vector field is symplectic (tangent to the family of Lagrangians) since each $K_\eps$ is Lagrangian. A vector field $X$ along a Lagrangian $L$ with $\iota_X \om$ closed is automatically Hamiltonian locally (since $H^1(L;\bR)$ classifies whether it's globally Hamiltonian). For $L$ simply connected (as will be the case for our local models near the vertex), $\iota_X \om$ is exact, giving a Hamiltonian $H$ with $\iota_{X}\om = -dH$.

More precisely: near each vertex $v_k$, the local isotopy $\Sigma_\eps^{(k)}$ is generated by a compactly supported Hamiltonian $H_k^\eps$ on $B_{r_k}(v_k)$ (the local smoothing is a Lagrangian isotopy in the ball, and Lagrangian isotopies in a simply-connected domain are Hamiltonian by the Weinstein Lagrangian isotopy theorem). These local Hamiltonians can be extended to zero outside the respective balls (using cutoff functions), giving a globally defined Hamiltonian $H_\eps = \sum_k \chi_k H_k^\eps$ that generates the isotopy from $K_\eps$ to $K_{\eps'}$ for $\eps, \eps' \in (0, \eps_0)$.

\begin{lemma}[Weinstein's Lagrangian isotopy theorem, local version]
\label{lem:weinstein-lag-isotopy}
Let $\{L_t\}_{t \in [0,1]}$ be a smooth isotopy of compact Lagrangian submanifolds in a symplectic manifold $(M, \om)$ (possibly with boundary, with $L_t$ fixed near the boundary). Then there exists a smooth family of Hamiltonian diffeomorphisms $\varphi_t$ with $\varphi_t(L_0) = L_t$.
\end{lemma}

\begin{proof}[Proof of Lemma \ref{lem:weinstein-lag-isotopy}]
The velocity vector field $\xi_t = \frac{d}{dt}L_t$ is a vector field along $L_t$ such that $\iota_{\xi_t}\om|_{L_t}$ is a closed 1-form (since $L_t$ is Lagrangian for all $t$, differentiating $\om|_{L_t} = 0$ in $t$ gives $\mathcal{L}_{\xi_t}\om|_{L_t} = 0$, which means $d(\iota_{\xi_t}\om)|_{L_t} = 0$). If $L_t$ is simply connected (or $H^1(L_t;\bR) = 0$), this 1-form is exact: $\iota_{\xi_t}\om = dH_t$ on $L_t$ for some function $H_t$ on $L_t$. Extending $H_t$ to a neighborhood by a cutoff function and then applying the Weinstein neighborhood theorem (which gives a symplectomorphism from a neighborhood of $L_t$ to a neighborhood of the zero section in $T^*L_t$), we extend $H_t$ to a globally defined function on $M$ (compactly supported near $L_t$). The resulting Hamiltonian vector field $X_{H_t}$ satisfies $X_{H_t}|_{L_t} = \xi_t$, hence the Hamiltonian flow $\varphi_t^{H}$ generated by $H_t$ satisfies $\varphi_t^H(L_0) = L_t$.

If $H^1(L_t;\bR) \neq 0$, then $\iota_{\xi_t}\om|_{L_t}$ need not be exact, and the isotopy is only symplectic (not necessarily Hamiltonian). However, for our local vertex models, $L_t$ is diffeomorphic to a disk (near the vertex, with boundary on the edges), hence simply connected, so the Hamiltonian condition is satisfied.
\end{proof}

Applying Lemma \ref{lem:weinstein-lag-isotopy} to each local isotopy $\{\Sigma_\eps^{(k)}\}$ and summing the resulting (compactly supported) Hamiltonians, we get a global Hamiltonian $H_\eps$ generating the isotopy of $K_\eps$. The isotopy is smooth in $\eps$ for $\eps \in (0, \eps_0)$.

Finally, \textbf{global reparametrization}: Let $\psi \colon (0,1] \to (0, \eps_0)$ be any smooth bijection (e.g., $\psi(t) = \eps_0 t/(1 + (1-\eps_0^{-1})(1-t))$). Setting $\tilde K_t = K_{\psi(t)}$ gives a Hamiltonian isotopy of smooth Lagrangian surfaces parameterized by $(0,1]$.
\end{proof}

\subsection{Continuity at $t = 0$}

\begin{proposition}[Continuity at $K_0 = K$]
\label{prop:continuity-at-zero}
As $\eps \to 0$, the surface $K_\eps$ converges to $K$ in the Hausdorff metric.
\end{proposition}

\begin{proof}
Outside the vertex balls $\bigcup_k B_{r_k/2}(v_k)$, we have $K_\eps = K$ for all $\eps$, so Hausdorff distance is controlled by the behavior near the vertices. Inside $B_{r_k/2}(v_k)$, the local smoothing $\Sigma_\eps^{(k)}$ converges to the local model $\Pi_{v_k} = K \cap B_{r_k/2}(v_k)$ in the Hausdorff metric as $\eps \to 0$, by Proposition \ref{prop:local-exist}(iii). Since there are finitely many vertices, the Hausdorff distance $d_H(K_\eps, K) \to 0$ as $\eps \to 0$.
\end{proof}

\section{Proof of the Main Theorem}
\label{sec:proof}

\begin{proof}[Proof of Theorem \ref{thm:main}]
Let $K$ be a polyhedral Lagrangian surface in $(\bR^4, \om)$ with exactly 4 faces meeting at every vertex.

\textbf{Step 1: Local normal form.} By Lemma \ref{lem:legendrian-link}, at each vertex $v$ the link $\Lk(v,K)$ is a Legendrian quadrilateral in $(S^3, \xi_{\rm std})$. By Lemma \ref{lem:normal-form}, the 4 Lagrangian planes at $v$ can be brought into the standard form by a linear symplectomorphism.

\textbf{Step 2: Local Lagrangian smoothings.} By Proposition \ref{prop:local-exist}, each vertex $v_k$ admits a local Lagrangian smoothing $\{\Sigma_\eps^{(k)}\}_{0<\eps<\eps_0}$ within a ball $B_{r_k}(v_k)$, degenerating to the original faces as $\eps\to 0$.

\textbf{Step 3: Global assembly.} Choose disjoint balls $B_k = B_{r_k}(v_k)$ around the vertices. Define $K_\eps$ by replacing $K \cap B_{r_k/2}(v_k)$ with $\Sigma_\eps^{(k)}$ for each $k$, and keeping $K$ elsewhere. By Lemma \ref{lem:edge-compat} and the construction, $K_\eps$ is a smooth embedded Lagrangian surface for each $\eps \in (0, \eps_0)$.

\textbf{Step 4: Hamiltonian isotopy.} By Proposition \ref{prop:hamiltonian-isotopy} and Lemma \ref{lem:weinstein-lag-isotopy}, the family $\{K_\eps\}$ is a Hamiltonian isotopy. After reparametrizing by $t \in (0,1]$, we get a smooth family $\{K_t\}_{t \in (0,1]}$ of smooth Lagrangian submanifolds with $K_t$ being Hamiltonian-isotopic to $K_1$ for all $t$.

\textbf{Step 5: Topological isotopy and boundary value.} By Proposition \ref{prop:continuity-at-zero}, the family extends continuously to $t = 0$ with $K_0 = K$ (in the Hausdorff metric, hence as a topological isotopy). The set-valued map $t \mapsto K_t$ is continuous on $[0,1]$.

Therefore $\{K_t\}_{t \in [0,1]}$ is a Lagrangian smoothing of $K$ in the sense of Definition 1.2, and $K$ admits a Lagrangian smoothing.
\end{proof}

\section{Remarks and Complements}

\begin{remark}[Why 4 faces and not more]
The condition that exactly 4 faces meet at every vertex is crucial. With more faces (e.g., 6 faces), the Legendrian link $\Lk(v,K)$ would be a Legendrian hexagon, and the cone over it might not admit a smooth Lagrangian filling. For 4 faces, the Legendrian link is always the standard Legendrian unknot (with $\mathrm{tb} = -1$), which is the unique Legendrian knot that bounds an exact Lagrangian disk in the standard Weinstein ball. This is the topological reason why 4 faces is the magic number for smooth Lagrangian filling.
\end{remark}

\begin{remark}[Comparison with the general case]
For a general polyhedral Lagrangian surface (with variable numbers of faces at vertices), the answer to Abouzaid's question would be: no, there need not exist a Lagrangian smoothing, because the Legendrian links at vertices with $\neq 4$ faces may obstruct filling. For instance, if 3 faces meet at a vertex, the Legendrian link is a Legendrian triangle, which may have $\mathrm{tb} \neq -1$ and fail to bound an exact Lagrangian disk.
\end{remark}

\begin{remark}[The Lagrangian cross]
The correct local model for the smoothing at a vertex is closely related to the \emph{Lagrangian cross} studied by Polterovich and Eliashberg. The smooth fiber of the map $(z_1, z_2) \mapsto z_1 z_2 \in \bC$ over a real value $\eps > 0$ gives a Lagrangian surface (the ``Milnor fiber'' of $z_1 z_2$) that smooths the union of coordinate planes. The condition that $z_1 z_2 \in \bR$ restricts to the real-Lagrangian setting, and the smoothing of the local model is exactly the deformation of the Lagrangian cross.
\end{remark}

\begin{remark}[Relation to Weinstein manifolds]
The proof crucially uses the fact that $\bR^4$ is a Weinstein manifold (or more precisely, a Liouville domain when restricted to a ball). The Weinstein neighborhood theorem for Lagrangian submanifolds (stating that a Lagrangian has a tubular neighborhood symplectomorphic to a neighborhood of the zero section in $T^*L$) is used both in the local Lagrangian filling construction and in extending Hamiltonians from local to global.
\end{remark}

\section{Summary}

We have proved Theorem \ref{thm:main}: every polyhedral Lagrangian surface $K \subset \bR^4$ with exactly 4 faces meeting at every vertex admits a Lagrangian smoothing. The proof has three main components:

\begin{enumerate}
  \item \textbf{Local analysis (Sections \ref{sec:vertex}--\ref{sec:local-smooth})}: The link of each vertex is a Legendrian quadrilateral in $(S^3, \xi_{\rm std})$, which is topologically a circle and admits an exact Lagrangian disk filling in the Weinstein 4-ball. This provides the local Lagrangian smoothing at each vertex.

  \item \textbf{Global assembly (Section \ref{sec:global})}: The local smoothings are combined using the Weinstein neighborhood theorem and compactly supported Hamiltonians to produce a globally smooth Lagrangian surface $K_\eps$ for each small $\eps > 0$.

  \item \textbf{Isotopy (Section \ref{sec:global})}: The family $K_\eps$ is a Hamiltonian isotopy (by Weinstein's Lagrangian isotopy theorem, applicable since the local surfaces are simply connected) that extends continuously to the singular surface $K_0 = K$.

\end{enumerate}

The answer to Abouzaid's question is \textbf{YES}: every polyhedral Lagrangian surface in $\bR^4$ with exactly 4 faces at every vertex has a Lagrangian smoothing.

\end{document}
